\chapter{Medians and Order Statistics}

Sorting, as you have now seen several times over, is one of the great
workhorses of algorithms. If you have a pile of numbers and you want to know
which is the largest, or the smallest, or the one sitting right in the middle,
the obvious move is to sort the whole pile and then read off the position you
care about. That works. But it is a little like alphabetizing your entire
bookshelf just to find one book. Surely, if all you want is \emph{one} item ---
the middle one, the third-largest, the value that splits the data in two --- you
should not have to put \emph{everything} in order first.

That nagging feeling is exactly right, and this chapter is about turning it into
a precise, fast algorithm. We are going to learn how to find the $k$-th smallest
element of a list in \emph{linear} time --- $\Th(n)$, the same order of growth as
simply \emph{looking} at every element once. That is striking. Reading all the
input already costs $\Th(n)$, so we will be doing the selection job for, in a
sense, the price of admission.

\begin{intuitionbox}
Sorting answers \emph{every} ``what is in position $i$?'' question at once, and
pays $\Th(n\log n)$ for the privilege. \textbf{Selection} answers just \emph{one}
such question. The whole chapter is the discovery that answering one question can
be done for $\Th(n)$ --- asymptotically free compared to sorting --- if you are
clever about throwing away the parts of the input that cannot possibly contain
the answer.
\end{intuitionbox}

\section{The selection problem}

Let us pin down what we are after. Imagine a list (an array) of $n$ numbers. They
are not sorted; they sit in whatever order they came. We will assume, to keep
things clean, that the numbers are \emph{distinct} --- no ties. (Ties can be
handled with a little care, but they cloud the ideas, so set them aside for now.)

\begin{definition}[Order statistic]\label{ch09:def-order-statistic}
The \textbf{$i$-th order statistic} of a set of $n$ numbers is the $i$-th smallest
of them. So the $1$st order statistic is the minimum, the $n$-th order statistic
is the maximum, and the $i$-th order statistic for a general $i$ is the value that
would land in position $i$ if you sorted the list into increasing order.
\end{definition}

\begin{definition}[Median]\label{ch09:def-median}
The \textbf{median} is the ``halfway point'' of the set. When $n$ is odd there is a
single middle element: the median is the $i$-th order statistic for
$i = (n+1)/2$. When $n$ is even there are \emph{two} middles; the
\textbf{lower median} is at $i = \floor{(n+1)/2} = n/2$ and the \textbf{upper
median} is at $i = \ceil{(n+1)/2} = n/2 + 1$. In this book, when we say ``the
median'' without qualification we mean the lower median, so that the formula
$i = \floor{(n+1)/2}$ always works.
\end{definition}

\begin{focusbox}
The \textbf{selection problem}: given an array $A$ of $n$ distinct numbers and an
integer $i$ with $1 \le i \le n$, return the element of $A$ that is larger than
exactly $i-1$ of the others --- the $i$-th smallest. The median is the special
case $i = \floor{(n+1)/2}$.
\end{focusbox}

Let us make the indices concrete, because this is where beginners stumble. Take
the array
\[
A = [\,7,\ 3,\ 9,\ 1,\ 5,\ 8,\ 2\,], \qquad n = 7.
\]
If we sorted it we would get $[\,1,2,3,5,7,8,9\,]$. So:
\begin{itemize}
  \item the $1$st order statistic (minimum) is $1$;
  \item the $7$th order statistic (maximum) is $9$;
  \item the $3$rd order statistic is $3$;
  \item the median, $i = \floor{(7+1)/2} = 4$, is $5$.
\end{itemize}
Notice the $i$-th order statistic is a \emph{value}, not a position in the original
array. In the unsorted $A$, the median value $5$ happens to live at index $5$, but
that coincidence means nothing. We care about its rank, not its address.

\begin{trapbox}
``$i$-th smallest'' is a \textbf{1-based rank}. The $1$st smallest is the minimum,
not the second-smallest, and not whatever sits at array index $1$. A huge fraction
of off-by-one bugs in selection code come from sliding between ``rank $i$'' and
``array index $i$.'' Keep them separate in your head: rank is a property of the
\emph{sorted order}; index is a property of \emph{where the value currently sits}.
\end{trapbox}

\subsection{The easy cases warm us up: minimum and maximum}

Before the general problem, look at the very simplest order statistics, $i=1$ and
$i=n$. Finding the minimum is the canonical first algorithm everyone writes: walk
through the array holding onto the smallest thing seen so far.

\begin{algorithm}[H]
\caption{\textsc{Minimum}$(A)$}\label{ch09:alg-minimum}
\KwIn{An array $A[1\,..\,n]$ of numbers, $n \ge 1$.}
\KwOut{The smallest element of $A$.}
$\textit{min} \gets A[1]$\;
\For{$j \gets 2$ \KwTo $n$}{
  \If{$A[j] < \textit{min}$}{
    $\textit{min} \gets A[j]$\;
  }
}
\Return $\textit{min}$\;
\end{algorithm}

This does exactly $n-1$ comparisons, one per element after the first. Can we do
better? No --- and there is a clean argument why. Think of each comparison as a
little tournament match between two values; the loser is the larger one (it cannot
be the minimum). To \emph{know} a particular element is the minimum, every other
element must have lost at least one match, because an element that never lost
might still be smaller than our claimed winner. There are $n-1$ ``other''
elements, each needing to lose at least once, and each comparison produces exactly
one loss. So you need at least $n-1$ comparisons. \textsc{Minimum} is optimal.

\begin{walkbox}
Suppose I ask for the minimum \emph{and} the maximum of the same array. The lazy
way runs \textsc{Minimum} and \textsc{Maximum} separately: $2(n-1)$ comparisons.
Can you do it in roughly $\tfrac{3}{2}n$ comparisons instead? Hint: process the
elements in \emph{pairs}. Compare the two members of a pair to each other first;
the smaller is a candidate for the min, the larger a candidate for the max. Count
how many comparisons that costs per pair, and you will see where the $\tfrac{3}{2}$
comes from.
\end{walkbox}

The minimum and maximum are gentle because $i$ is fixed at an extreme. The
interesting trouble is the \emph{general} $i$, and in particular the median, which
sits as far from both ends as possible. That is what the rest of the chapter
attacks.

\section{Selection by sorting: the baseline to beat}

The honest first algorithm is the one we already grumbled about: sort, then index.

\begin{algorithm}[H]
\caption{\textsc{SelectBySorting}$(A, i)$}\label{ch09:alg-sortselect}
\KwIn{An array $A[1\,..\,n]$ of distinct numbers; a rank $i$, $1 \le i \le n$.}
\KwOut{The $i$-th smallest element of $A$.}
$B \gets \textsc{Sort}(A)$ \tcp*{$B$ is $A$ in increasing order}
\Return $B[i]$\;
\end{algorithm}

It is correct by definition: the $i$-th order statistic is, literally, the element
in position $i$ of the sorted array. With a good comparison sort like merge sort or
heapsort, it costs $\Th(n\log n)$. That is our baseline. Everything we build from
here is an attempt to beat it --- to get the single answer without paying for the
full ordering.

\begin{aestheticbox}
Why fight for the difference between $n\log n$ and $n$? Because selection is a
\emph{subroutine}, called over and over inside bigger jobs. Order statistics power
the ``median-of-medians'' pivots that make quicksort robust; they let databases
answer ``the 95th percentile latency'' without sorting a billion rows; they sit
inside $k$-nearest-neighbor search, robust statistics (the median is far more
outlier-resistant than the mean), and load balancers picking a cutoff. When a
routine is called $n$ times, shaving a $\log n$ factor off each call is the
difference between a system that scales and one that melts.
\end{aestheticbox}

\section{Quickselect: partition, then recurse on one side}

Here is the key shift in thinking. In Chapter~7 you met \textsc{Partition}, the
engine inside quicksort. It picks a \emph{pivot}, then rearranges the array so that
everything smaller than the pivot sits to its left and everything larger sits to
its right, with the pivot itself landing in its final sorted position. Quicksort
then recurses on \emph{both} sides.

But for selection we do not need both sides. After partitioning, the pivot is
sitting at its true sorted rank. If that rank happens to be exactly $i$, we are
done --- we found it. If the pivot's rank is bigger than $i$, the answer must lie
to the \emph{left} of the pivot, so we recurse there and forget the right half
entirely. If the pivot's rank is smaller than $i$, the answer is to the
\emph{right}, and we recurse there, forgetting the left.

\begin{intuitionbox}
Quicksort recurses on both halves because it wants \emph{everything} sorted.
Selection recurses on \emph{one} half because it wants \emph{one} element. Throwing
away half the array on each step is the whole trick: it turns the
$T(n) = 2T(n/2) + \Th(n)$ recurrence of quicksort (which gives $n\log n$) into
something closer to $T(n) = T(n/2) + \Th(n)$ (which gives $\Th(n)$).
\end{intuitionbox}

\subsection{Recalling \textsc{Partition}}

Let me restate the partition routine so this chapter stands on its own. We use the
Lomuto scheme: take the last element as the pivot, sweep a finger across the
array, and keep a wall such that everything left of the wall is $\le$ pivot.

\begin{algorithm}[H]
\caption{\textsc{Partition}$(A, p, r)$}\label{ch09:alg-partition}
\KwIn{Subarray $A[p\,..\,r]$.}
\KwOut{Rearranges $A[p\,..\,r]$ around the pivot $A[r]$; returns the pivot's final index $q$.}
$x \gets A[r]$ \tcp*{the pivot value}
$j \gets p - 1$ \tcp*{$j$ is the wall: $A[p\,..\,j]$ are all $\le x$}
\For{$k \gets p$ \KwTo $r-1$}{
  \If{$A[k] \le x$}{
    $j \gets j + 1$\;
    swap $A[j] \leftrightarrow A[k]$\;
  }
}
swap $A[j+1] \leftrightarrow A[r]$ \tcp*{drop the pivot just past the wall}
\Return $j + 1$\;
\end{algorithm}

After \textsc{Partition} returns $q$, three facts hold: $A[q]$ is the original
pivot; every element in $A[p\,..\,q-1]$ is less than $A[q]$; and every element in
$A[q+1\,..\,r]$ is greater than $A[q]$. So $A[q]$ is in its final sorted slot
\emph{within this subarray}. Its rank \emph{within the subarray} is $q - p + 1$
(the count of elements at positions $p$ through $q$).
Figure~\ref{ch09:fig-partition-state} pictures the three regions and how selection
chooses which one to keep.

\begin{figure}[H]
\centering
\begin{tikzpicture}[every node/.style={font=\small}]
  % three regions of the partitioned subarray
  \draw[targetgreen, thick, fill=targetgreen!10] (0,0) rectangle (4,1);
  \draw[trapred, thick, fill=trapred!15] (4,0) rectangle (5,1);
  \draw[sourceblue, thick, fill=sourceblue!10] (5,0) rectangle (9,1);
  \node[targetgreen] at (2,0.5) {all $< $ pivot};
  \node[trapred] at (4.5,0.5) {\footnotesize pivot};
  \node[sourceblue] at (7,0.5) {all $>$ pivot};
  % position labels
  \node[below] at (0,0) {\footnotesize $p$};
  \node[below] at (4,-0.02) {};
  \node[below=1pt] at (4.5,0) {\footnotesize $q$};
  \node[below] at (9,0) {\footnotesize $r$};
  % rank annotation
  \node[trapred, above] at (4.5,1.1) {\footnotesize rank $k = q-p+1$};
  % decisions
  \draw[->, targetgreen, thick] (2,-0.7) -- (2,-0.15);
  \node[targetgreen, below=22pt, align=center] at (2,0) {\footnotesize if $i<k$:\\[-2pt]\footnotesize recurse here, rank $i$};
  \draw[->, sourceblue, thick] (7,-0.7) -- (7,-0.15);
  \node[sourceblue, below=22pt, align=center] at (7,0) {\footnotesize if $i>k$:\\[-2pt]\footnotesize recurse here, rank $i-k$};
  \draw[->, trapred, thick] (4.5,1.9) -- (4.5,1.25);
  \node[trapred, above=30pt, align=center] at (4.5,0) {\footnotesize if $i=k$: done};
\end{tikzpicture}
\caption{After \textsc{Partition}, the pivot sits at its true rank $k=q-p+1$.
Selection keeps only the one shaded region that can hold rank $i$ and discards the
rest --- including the pivot's whole other side.}
\label{ch09:fig-partition-state}
\end{figure}

\begin{trapbox}
The rank of the pivot \emph{inside the current subarray} is $k = q - p + 1$, not
$q$ itself. When you recurse, you must convert between ``rank within the whole
problem'' and ``rank within the subarray.'' Forgetting the $-p+1$ shift, or
forgetting to subtract $k$ when you recurse right, is the number-one bug in
quickselect. We will walk a numeric trace precisely to nail this down.
\end{trapbox}

\subsection{Randomized quickselect}

We will use a \emph{randomized} pivot: before partitioning, swap a uniformly random
element into the last slot. This is the same idea you saw stabilize quicksort ---
it defends against adversarial inputs and makes the analysis about randomness in
\emph{our} coin flips rather than luck in the input.

\begin{algorithm}[H]
\caption{\textsc{RandomizedPartition}$(A, p, r)$}\label{ch09:alg-rpartition}
\KwIn{Subarray $A[p\,..\,r]$.}
\KwOut{A partition around a randomly chosen pivot; returns the pivot's index.}
$t \gets \textsc{Random}(p, r)$ \tcp*{uniform integer in $[p, r]$}
swap $A[r] \leftrightarrow A[t]$\;
\Return \textsc{Partition}$(A, p, r)$\;
\end{algorithm}

\begin{algorithm}[H]
\caption{\textsc{RandomizedSelect}$(A, p, r, i)$}\label{ch09:alg-rselect}
\KwIn{Subarray $A[p\,..\,r]$; rank $i$ with $1 \le i \le r-p+1$, the rank \emph{within this subarray}.}
\KwOut{The $i$-th smallest element of $A[p\,..\,r]$.}
\If{$p = r$}{
  \Return $A[p]$ \tcp*{one element, must be the answer}
}
$q \gets \textsc{RandomizedPartition}(A, p, r)$\;
$k \gets q - p + 1$ \tcp*{rank of the pivot within $A[p\,..\,r]$}
\uIf{$i = k$}{
  \Return $A[q]$ \tcp*{the pivot \emph{is} the answer}
}
\uElseIf{$i < k$}{
  \Return \textsc{RandomizedSelect}$(A, p, q-1, i)$ \tcp*{answer is on the left}
}
\Else{
  \Return \textsc{RandomizedSelect}$(A, q+1, r, i-k)$ \tcp*{answer on the right; shift the rank}
}
\end{algorithm}

Stare at the last line, because that shift is the subtle part. When we recurse
\emph{right}, the new subarray $A[q+1\,..\,r]$ has had $k$ elements chopped off its
front (positions $p$ through $q$). An element that was the $i$-th smallest in the
old subarray is now the $(i-k)$-th smallest in the new one, because $k$ smaller
elements are no longer in the picture. That is why we pass $i - k$, not $i$. When
we recurse \emph{left}, no elements smaller than the target were removed, so the
rank $i$ is unchanged.

\begin{focusbox}
Call \textsc{RandomizedSelect}$(A, 1, n, i)$ to find the $i$-th smallest of the
whole array. The procedure never sorts; it only partitions and then dives into the
\emph{one} side that can contain the answer, adjusting the rank as it goes.
\end{focusbox}

\subsection{A full numeric trace}

Let us run it on
\[
A = [\,7,\ 3,\ 9,\ 1,\ 5,\ 8,\ 2\,], \qquad n = 7,
\]
asking for the median, $i = 4$. To make the trace reproducible I will pretend the
random pivot choices come out the way I describe; the point is to see the
\emph{logic}, not the dice.

\textbf{Step 1.} Call \textsc{RandomizedSelect}$(A, 1, 7, 4)$. Say the random pivot
is $A[5] = 5$. Swap it to the end: $A = [7,3,9,1,2,8,5]$. Now partition around
$x=5$.
Sweeping $k = 1\ldots 6$: $7 > 5$ skip; $3 \le 5$ advance wall, $A$ becomes
$[3,7,9,1,2,8,5]$; $9>5$ skip; $1\le 5$ advance wall, $[3,1,9,7,2,8,5]$; $2 \le 5$
advance wall, $[3,1,2,7,9,8,5]$; $8>5$ skip. Finally drop the pivot just past the
wall (position $4$): $A = [3,1,2,5,9,8,7]$, and $q = 4$.

The pivot $5$ sits at index $q=4$. Its rank within $A[1\,..\,7]$ is
$k = q - p + 1 = 4 - 1 + 1 = 4$. We wanted $i = 4$. Since $i = k$, \textbf{we are
done}: the median is $5$. One partition, no recursion at all. That was lucky ---
the random pivot happened to be the median itself. Let us run a less lucky example
to see the recursion fire.

\textbf{A second trace, $i = 6$ (the $6$th smallest).} Same starting
$A = [7,3,9,1,5,8,2]$, call \textsc{RandomizedSelect}$(A,1,7,6)$. Suppose the random
pivot is $A[2] = 3$. Swap to end: $A = [7,2,9,1,5,8,3]$, partition around $x=3$.
Sweep: $7>3$ skip; $2\le 3$ wall to index 1, $[2,7,9,1,5,8,3]$; $9>3$ skip;
$1\le3$ wall to index 2, $[2,1,9,7,5,8,3]$; $5>3$ skip; $8>3$ skip. Drop pivot at
index $3$: $A=[2,1,3,7,5,8,9]$, $q = 3$. Rank $k = 3 - 1 + 1 = 3$.

We want $i = 6$; since $6 > 3$, recurse right on $A[4\,..\,7] = [7,5,8,9]$ with new
rank $i - k = 6 - 3 = 3$. Call \textsc{RandomizedSelect}$(A, 4, 7, 3)$. Suppose the
pivot is $A[6] = 8$. Swap to end ($A[7]$): $A[4\,..\,7]$ becomes $[7,5,9,8]$,
partition around $x=8$. Sweep over $[7,5,9]$: $7\le 8$ wall, then $5 \le 8$ wall,
then $9 > 8$ skip; drop pivot: $A[4\,..\,7] = [7,5,8,9]$, $q = 6$. Rank within this
subarray is $k = 6 - 4 + 1 = 3$. We wanted $i = 3$ here, and $i = k = 3$, so
\textbf{return} $A[6] = 8$. And indeed the sorted array is $[1,2,3,5,7,8,9]$, whose
$6$th smallest is $8$. Correct.

\begin{labbox}
Trace \textsc{RandomizedSelect}$([7,3,9,1,5,8,2], 1, 7, 2)$ by hand, choosing the
pivots however you like, and confirm you reach $2$ (the $2$nd smallest). Do it a
second time with \emph{different} pivot choices and confirm you still reach $2$.
The pivots change the \emph{path}; the answer never moves. That invariance --- same
output regardless of the coin flips --- is exactly what makes a randomized
algorithm trustworthy.
\end{labbox}

\subsection{Why quickselect is correct}

Correctness rests on a single invariant, proved by induction on the size of the
subarray $n' = r - p + 1$.

\begin{formalbox}
\textbf{Claim.} \textsc{RandomizedSelect}$(A, p, r, i)$, called with
$1 \le i \le r - p + 1$, returns the $i$-th smallest element of $A[p\,..\,r]$.

\emph{Base case.} If $p = r$ the subarray has one element, $n' = 1$, and the only
legal rank is $i = 1$. The procedure returns $A[p]$, which is indeed the $1$st (and
only) smallest. True.

\emph{Inductive step.} Assume the claim for all subarrays smaller than $A[p\,..\,r]$.
After \textsc{RandomizedPartition} returns $q$, the pivot $A[q]$ is greater than
exactly the $k-1$ elements in $A[p\,..\,q-1]$ and less than the elements in
$A[q+1\,..\,r]$, so the pivot's true rank in $A[p\,..\,r]$ is $k = q-p+1$.
\begin{itemize}
  \item If $i = k$, the $i$-th smallest is precisely the pivot; we return $A[q]$.
  \item If $i < k$, the $i$-th smallest of $A[p\,..\,r]$ lies among the $k-1$
    elements strictly below the pivot, i.e.\ in $A[p\,..\,q-1]$, and at the
    \emph{same} rank $i$ (no smaller element was removed). That subarray is
    smaller, so by induction the recursive call returns the right value.
  \item If $i > k$, the target lies in $A[q+1\,..\,r]$. Exactly $k$ elements (the
    pivot and the $k-1$ below it) were removed from the front, all of them smaller
    than the target, so the target's rank in the new subarray is $i - k$. Again by
    induction the recursive call is correct.
\end{itemize}
In all three cases the value returned is the $i$-th smallest of $A[p\,..\,r]$.
\qed
\end{formalbox}

\section{How fast is quickselect?}

Now the heart of the matter. Each call does $\Th(n)$ work to partition, then
recurses on \emph{one} side. The size of that side depends on where the pivot fell
--- and that is random. So we get a range of behaviors.

\subsection{The worst case is bad}

Imagine the pivot is, every single time, the smallest remaining element. Then the
partition peels off just one element and recurses on the other $n-1$. The sizes go
$n, n-1, n-2, \ldots, 1$, and the total work is
\[
\Th(n) + \Th(n-1) + \cdots + \Th(1) = \Th\!\left(\sum_{j=1}^{n} j\right) = \Th(n^2).
\]
So in the worst case, quickselect is \emph{quadratic} --- no better than the naive
bubble-style approaches. With random pivots this worst case is astronomically
unlikely (it requires the random choice to land on the extreme every time), but it
is possible, and a worst-case-conscious engineer will not be satisfied. We will fix
it in the next section. First, let us see why the \emph{expected} case is excellent.

\subsection{Expected linear time: the intuition first}

\begin{intuitionbox}
A random pivot splits the array somewhere. On average it lands near the middle ---
not exactly at the median, but ``balanced enough.'' A pivot is what we will call
\textbf{good} if it falls in the middle half of the values, i.e.\ its rank is
between $n/4$ and $3n/4$. The chance of that is $1/2$ (half the positions qualify).
A good pivot guarantees the side we recurse into has at most $3/4$ of the elements.
So \emph{at least half the time}, the problem shrinks to at most three-quarters its
size. Shrinking by a constant factor a constant fraction of the time is exactly the
condition for a geometric series --- and a geometric series of linear-work terms
sums to linear total work.
\end{intuitionbox}

Let us make that intuition into a clean bound, then sketch the sharper expectation
argument.

\paragraph{A clean bound via ``good'' pivots.}
Call a partition step \emph{lucky} if the pivot's rank lands in the middle half,
between $\ceil{n/4}$ and $\floor{3n/4}$. There are about $n/2$ such ranks out of
$n$, so a uniformly random pivot is lucky with probability $\ge 1/2$. When a step
is lucky, whichever side we recurse on has size at most $3n/4$ (because at least
$n/4$ elements sit on each side of a middle-half pivot). Now think of the sequence
of subproblem sizes. Each step does $\le c\,n$ work for the current $n$. On average
we need at most $2$ tries to get a lucky step (a geometric random variable with
success probability $1/2$ has mean $2$). So, in expectation, the size shrinks by a
factor of $3/4$ every $2$ steps. Writing $T(n)$ for the expected work,
\[
T(n) \le T\!\left(\tfrac{3}{4}n\right) + (\text{expected work to reach a lucky step}).
\]
The expected work to reach a lucky step is at most $2 c n$ (two partitions, each
$\le cn$, since sizes only shrink). So
\[
T(n) \le T\!\left(\tfrac{3}{4}n\right) + 2 c n.
\]
Unrolling this geometric recurrence,
\[
T(n) \le 2cn\left(1 + \tfrac{3}{4} + \left(\tfrac{3}{4}\right)^2 + \cdots\right)
      = 2cn \cdot \frac{1}{1 - 3/4} = 2cn \cdot 4 = 8cn = \Th(n).
\]
The geometric series $\sum (3/4)^j$ converges to $4$, a \emph{constant}, so the
whole thing is linear.

\begin{focusbox}
The reason quickselect is linear and quicksort is $n\log n$ comes down to one line:
selection recurses on \emph{one} side, so its sizes form a \emph{geometric}
(shrinking) sequence that sums to $\Th(n)$; sorting recurses on \emph{both} sides,
so it must touch every level fully, paying $\Th(n)$ per level across $\Th(\log n)$
levels.
\end{focusbox}

Figure~\ref{ch09:fig-geometric} makes the contrast visual: quicksort's tree is full
and fat, while quickselect walks a single shrinking path down one branch.

\begin{figure}[H]
\centering
\begin{tikzpicture}[every node/.style={font=\footnotesize}, node distance=4mm]
  % quickselect: a single chain of shrinking bars
  \begin{scope}
    \node at (1.5,3.2) {\textbf{Quickselect: one shrinking path}};
    \draw[sourceblue, fill=sourceblue!20] (0,2.3) rectangle (3,2.7);
    \node at (1.5,2.5) {$n$};
    \draw[sourceblue, fill=sourceblue!20] (0,1.6) rectangle (2.25,2.0);
    \node at (1.125,1.8) {$\le \tfrac34 n$};
    \draw[sourceblue, fill=sourceblue!20] (0,0.9) rectangle (1.69,1.3);
    \node at (0.84,1.1) {$\le (\tfrac34)^2 n$};
    \draw[sourceblue, fill=sourceblue!20] (0,0.2) rectangle (1.27,0.6);
    \node at (0.63,0.4) {$\cdots$};
    \node[align=center] at (1.5,-0.4) {sizes sum to a\\geometric series $=\Th(n)$};
  \end{scope}
  % quicksort: full binary tree
  \begin{scope}[xshift=6.5cm]
    \node at (1.5,3.2) {\textbf{Quicksort: every level full}};
    \node (a) at (1.5,2.5) {$n$};
    \node (b) at (0.5,1.5) {$\tfrac{n}{2}$};
    \node (c) at (2.5,1.5) {$\tfrac{n}{2}$};
    \node (d) at (0,0.5) {$\tfrac{n}{4}$};
    \node (e) at (1,0.5) {$\tfrac{n}{4}$};
    \node (f) at (2,0.5) {$\tfrac{n}{4}$};
    \node (g) at (3,0.5) {$\tfrac{n}{4}$};
    \draw (a)--(b) (a)--(c) (b)--(d) (b)--(e) (c)--(f) (c)--(g);
    \node[align=center] at (1.5,-0.4) {$\Th(n)$ per level,\\$\Th(\log n)$ levels $=\Th(n\log n)$};
  \end{scope}
\end{tikzpicture}
\caption{Selection descends a single branch whose sizes shrink by a constant factor,
so the total work is a convergent geometric series and stays linear. Sorting must
fill out the entire tree, paying linear work on each of its $\Th(\log n)$ levels.}
\label{ch09:fig-geometric}
\end{figure}

\subsection{Expected linear time: the sharper analysis}

The ``good pivot'' bound is honest and gives the right answer, $\Th(n)$. CLRS
derives the constant more precisely with indicator random variables; let us walk
that path gently, because it is a beautiful use of the tools from
Chapter~5.

When we partition a subarray of size $n$, the pivot is equally likely to have any
rank from $1$ to $n$, each with probability $1/n$. If the pivot rank is $k$, the
recursion goes left into a subarray of size $k-1$ or right into a subarray of size
$n-k$ (or stops, if we got lucky and $i = k$). To get an \emph{upper} bound on the
expected time, assume the adversary always sends us into the \emph{larger} of the
two sides --- that can only slow us down. The larger side has size
$\max(k-1,\,n-k)$.

Let $T(n)$ be the maximum, over all $i$, of the expected running time on an array
of $n$ elements. Then
\[
T(n) \le \frac{1}{n}\sum_{k=1}^{n} T\big(\max(k-1,\,n-k)\big) + \Th(n).
\]
As $k$ ranges over $1,\dots,n$, the quantity $\max(k-1, n-k)$ runs through the
values $\floor{n/2}, \dots, n-1$, and each such size is hit \emph{at most twice}.
So
\[
T(n) \le \frac{2}{n}\sum_{m = \floor{n/2}}^{n-1} T(m) + \Th(n).
\]
Now guess $T(n) \le c\,n$ for a constant $c$ to be chosen, and substitute:
\[
T(n) \le \frac{2}{n}\sum_{m=\floor{n/2}}^{n-1} c\,m + \Th(n)
      = \frac{2c}{n}\sum_{m=\floor{n/2}}^{n-1} m + \Th(n).
\]
The inner sum is a stretch of consecutive integers. Bounding it from above by the
full sum minus the small head,
\[
\sum_{m=\floor{n/2}}^{n-1} m = \sum_{m=1}^{n-1} m - \sum_{m=1}^{\floor{n/2}-1} m
\le \frac{(n-1)n}{2} - \frac{(n/2-1)(n/2)}{2}
\le \frac{3n^2}{8}.
\]
(You can verify the $3n^2/8$ by a little algebra; it is the cleanest constant that
survives the floors.) Plugging back in,
\[
T(n) \le \frac{2c}{n}\cdot\frac{3n^2}{8} + \Th(n) = \frac{3cn}{4} + \Th(n).
\]
Write the $\Th(n)$ term as $a n$ for a constant $a$. Then
\[
T(n) \le \frac{3cn}{4} + a n = c n - \left(\frac{cn}{4} - a n\right).
\]
If we choose $c \ge 4a$, the parenthesized quantity is $\ge 0$, so
$T(n) \le cn$, completing the induction. Hence $T(n) = \Th(n)$: quickselect runs in
\textbf{expected linear time}.

\begin{formalbox}
\textbf{Theorem.} On an array of $n$ distinct numbers, \textsc{RandomizedSelect}
runs in expected time $\Th(n)$, where the expectation is over the algorithm's own
random pivot choices and holds for every input and every rank $i$.
\end{formalbox}

\begin{complexitybox}
\textbf{Randomized quickselect.}
\begin{itemize}
  \item Expected time: $\Th(n)$ --- linear, for \emph{any} input.
  \item Worst-case time: $\Th(n^2)$ --- but only on a vanishingly unlikely sequence
    of pivot choices, not on any fixed input.
  \item Extra space: $\Th(1)$ if you turn the tail recursion into a loop;
    $\Th(\log n)$ expected stack depth otherwise. The array is permuted in place.
  \item Comparisons in the expected case: about $3.4n$ for the median with a careful
    implementation --- genuinely cheaper than sorting.
\end{itemize}
\end{complexitybox}

\begin{trapbox}
``Expected linear'' is \emph{not} ``always linear.'' For most applications the
expectation, plus the fact that bad luck is exponentially rare, is plenty. But if
you are building something where a rare quadratic blowup is unacceptable --- a
real-time system, or code facing an adversary who can see your random seed --- you
need a \emph{worst-case} guarantee. That is precisely what the next section
delivers.
\end{trapbox}

\section{Median of medians: worst-case linear selection}

Now for one of the most celebrated results in algorithms: you can find the $i$-th
order statistic in $\Th(n)$ time \emph{in the worst case}, with no randomness at
all. This is the 1973 algorithm of Blum, Floyd, Pratt, Rivest, and Tarjan, often
called \textsc{Select} or simply ``median of medians.''

The problem with quickselect's worst case was bad pivots. A random pivot is good
\emph{on average}, but a single unlucky run can pick the extreme every time. The
fix is to stop gambling and instead \emph{compute} a pivot that is guaranteed to be
reasonably central --- guaranteed to throw away a constant fraction of the array no
matter what. The genius is doing that guarantee-computation cheaply enough that it
does not wreck the linear budget.

\begin{intuitionbox}
We cannot afford to find the \emph{true} median to use as a pivot --- that is the
very problem we are trying to solve, so it would be circular. But we do not need the
true median. We need a pivot that is merely ``okay'' --- guaranteed to split off at
least, say, $30\%$ of the array on each side. The median-of-medians trick
manufactures such an okay pivot by finding the median of a small \emph{sample} of
medians, cleverly chosen so its quality can be proven.
\end{intuitionbox}

\subsection{The five steps}

The algorithm \textsc{Select}$(A, i)$ does the following.

\begin{enumerate}
  \item \textbf{Group.} Divide the $n$ elements into groups of $5$ (the last group
    may have fewer). There are $\ceil{n/5}$ groups.
  \item \textbf{Sort each group.} Find the median of each group of $5$. Since each
    group is tiny, you can sort it with a fixed number of comparisons --- a
    constant amount of work per group, $\Th(n)$ total.
  \item \textbf{Recurse for the pivot.} Collect the $\ceil{n/5}$ group medians into
    a list and \emph{recursively call \textsc{Select}} to find \emph{their} median.
    Call this value $x$, the ``median of medians.''
  \item \textbf{Partition around $x$.} Use $x$ as the pivot in a partition of the
    whole array. Let $k$ be $x$'s rank in $A$.
  \item \textbf{Recurse into one side.} Exactly as in quickselect: if $i = k$
    return $x$; if $i < k$ recurse left; if $i > k$ recurse right with rank
    $i - k$.
\end{enumerate}

Steps 4 and 5 are pure quickselect. The new ideas are steps 1--3, which replace
``pick a random pivot'' with ``compute a provably central pivot.''

\begin{algorithm}[H]
\caption{\textsc{Select}$(A, p, r, i)$ -- deterministic, worst-case linear}\label{ch09:alg-select}
\KwIn{Subarray $A[p\,..\,r]$ of distinct numbers; rank $i$, $1 \le i \le r-p+1$.}
\KwOut{The $i$-th smallest element of $A[p\,..\,r]$.}
$n \gets r - p + 1$\;
\If{$n \le 5$}{
  sort $A[p\,..\,r]$ directly and \Return $A[p + i - 1]$\;
}
divide $A[p\,..\,r]$ into $\ceil{n/5}$ groups of $5$ (last group may be smaller)\;
\For{each group}{
  sort it and record its median\;
}
let $M$ be the array of the $\ceil{n/5}$ group medians\;
$x \gets \textsc{Select}\big(M,\, 1,\, \abs{M},\, \floor{(\abs{M}+1)/2}\big)$ \tcp*{median of medians}
$q \gets \textsc{PartitionAround}(A, p, r, x)$ \tcp*{partition $A[p\,..\,r]$ using value $x$ as pivot}
$k \gets q - p + 1$\;
\uIf{$i = k$}{\Return $x$\;}
\uElseIf{$i < k$}{\Return \textsc{Select}$(A, p, q-1, i)$\;}
\Else{\Return \textsc{Select}$(A, q+1, r, i-k)$\;}
\end{algorithm}

Here \textsc{PartitionAround}$(A,p,r,x)$ is just \textsc{Partition} modified to use
a specified pivot \emph{value} $x$ (find $x$ in the subarray, swap it to the end,
and run the usual Lomuto sweep). Everything else mirrors quickselect.

\subsection{Why the pivot is guaranteed good}

This is the crux, and it is best seen as a picture. Lay out the $\ceil{n/5}$ groups
as columns, each column sorted from smallest (bottom) to largest (top), and arrange
the columns left to right by their median. The median-of-medians $x$ sits in the
middle column, at the middle height. Figure~\ref{ch09:fig-mom} draws a small case.

\begin{figure}[H]
\centering
\begin{tikzpicture}[scale=0.62, every node/.style={font=\small}]
  % draw a 5x7 grid of dots representing 7 groups of 5, sorted within columns
  \def\cols{7}
  \def\rows{5}
  \foreach \c in {1,...,7} {
    \foreach \rw in {1,...,5} {
      \fill[gray!55] (\c, \rw) circle (3pt);
    }
  }
  % highlight the median of each column (middle row)
  \foreach \c in {1,...,7} {
    \fill[quotientorange] (\c, 3) circle (4pt);
  }
  % the median of medians sits in the middle column (col 4), middle row
  \fill[trapred] (4, 3) circle (5.5pt);
  \node[trapred, below=10pt] at (4, 0.7) {$x$ (median of medians)};
  % shade the guaranteed-smaller region: lower-left block
  \draw[targetgreen, thick, rounded corners]
       (0.6, 0.6) rectangle (4.4, 3.4);
  \node[targetgreen] at (2.5, 4.5) {all $\le x$ (provably)};
  % shade the guaranteed-larger region: upper-right block
  \draw[sourceblue, thick, rounded corners]
       (3.6, 2.6) rectangle (7.4, 5.4);
  \node[sourceblue] at (6.4, 1.6) {all $\ge x$ (provably)};
  % axis hints
  \node[rotate=90] at (-0.1, 3) {\footnotesize larger within group $\uparrow$};
  \node at (4, -0.4) {\footnotesize columns sorted by their median $\rightarrow$};
\end{tikzpicture}
\caption{Seven groups of five, each column sorted upward, columns ordered left to
right by their median (orange). The median-of-medians $x$ (red) dominates the entire
green block ($\ge \tfrac{3}{10}n$ elements that are $\le x$) and is dominated by the
blue block (another $\ge \tfrac{3}{10}n$ elements that are $\ge x$).}
\label{ch09:fig-mom}
\end{figure}

Read the picture carefully. Each column is a group of $5$, sorted so larger
elements are higher. The orange dots are the group medians; columns are arranged so
their medians increase left to right. The red dot $x$ is the median of those
medians.

Now look at the green block: the lower-left region. Take any column \emph{to the
left} of $x$'s column (its median is $\le x$). In such a column, the median and the
two elements below it --- that is $3$ of the $5$ --- are all $\le$ the column's
median, hence $\le x$. The same holds for $x$'s own column. So in roughly half the
columns, at least $3$ elements are guaranteed $\le x$. By a symmetric argument
(the blue block), in roughly half the columns at least $3$ elements are guaranteed
$\ge x$.

\begin{formalbox}
\textbf{How many elements are guaranteed $\le x$?} At least half of the
$\ceil{n/5}$ groups contribute $3$ elements each that are $\le x$. Being careful
with the group containing $x$ and a possibly-undersized last group, the count of
elements guaranteed $\le x$ is at least
\[
3\left(\ceil{\tfrac{1}{2}\ceil{\tfrac{n}{5}}} - 2\right) \ge \frac{3n}{10} - 6.
\]
By symmetry, at least $\tfrac{3n}{10} - 6$ elements are guaranteed $\ge x$.
Therefore the side we recurse into in step 5 has size at most
\[
n - \left(\frac{3n}{10} - 6\right) = \frac{7n}{10} + 6.
\]
\end{formalbox}

\begin{focusbox}
The payoff in one line: the pivot $x$ is \emph{guaranteed} to throw away at least
about $\tfrac{3}{10}n$ elements, so the recursion never faces a subproblem bigger
than $\tfrac{7}{10}n + 6$. No randomness, no luck --- a worst-case promise. (Why
groups of $5$? See the box after the analysis.)
\end{focusbox}

\subsection{The recurrence, and why $5$ is magic}

Now we tally the work. Let $T(n)$ be the worst-case running time of \textsc{Select}
on $n$ elements.
\begin{itemize}
  \item Grouping into fives and finding each group's median: $\Th(n)$, since each of
    the $\ceil{n/5}$ groups costs a constant (sort $5$ items in $\le 7$
    comparisons).
  \item The recursive call to find the median of medians runs on $\ceil{n/5}$
    elements: $T(\ceil{n/5})$.
  \item Partitioning around $x$: $\Th(n)$.
  \item The recursive call in step 5 runs on at most $\tfrac{7n}{10} + 6$ elements:
    $T(\tfrac{7n}{10}+6)$.
\end{itemize}
Putting it together, for $n$ above some small threshold,
\[
T(n) \le T\!\left(\ceil{\tfrac{n}{5}}\right) + T\!\left(\tfrac{7n}{10} + 6\right) + \Th(n).
\]
This is an unusual recurrence: it has \emph{two} recursive calls, but on
subproblems whose sizes add up to \emph{less than} $n$. That is the whole game.
Notice $\tfrac{1}{5} + \tfrac{7}{10} = \tfrac{2}{10} + \tfrac{7}{10} = \tfrac{9}{10} < 1$.
The two children together are only $9/10$ of the parent, so the work per level
\emph{shrinks geometrically}, and a geometric series of linear terms is linear.

Let us verify by the substitution method. Guess $T(n) \le c\,n$ for all $n$ past
the threshold, with $c$ to be fixed and the $\Th(n)$ term written as $a n$:
\[
T(n) \le c\ceil{\tfrac{n}{5}} + c\left(\tfrac{7n}{10}+6\right) + a n
      \le \frac{cn}{5} + c + \frac{7cn}{10} + 6c + a n.
\]
Combine the $n$ terms: $\tfrac{cn}{5} + \tfrac{7cn}{10} = \tfrac{2cn + 7cn}{10} = \tfrac{9cn}{10}$.
So
\[
T(n) \le \frac{9cn}{10} + 7c + a n
      = cn - \left(\frac{cn}{10} - 7c - a n\right)
      = cn - \left(\Big(\tfrac{c}{10} - a\Big)n - 7c\right).
\]
We need the bracket $\ge 0$, i.e.\ $(\tfrac{c}{10} - a)n \ge 7c$. Choose
$c = 20a$, say. Then $\tfrac{c}{10} - a = 2a - a = a$, and the condition becomes
$a n \ge 7c = 140 a$, i.e.\ $n \ge 140$. So for all $n \ge 140$ (handle smaller $n$
by the brute-force base case), $T(n) \le cn = 20an$. Hence
\[
\boxed{T(n) = \Th(n)} \qquad \text{worst case.}
\]

\begin{formalbox}
\textbf{Theorem (Blum--Floyd--Pratt--Rivest--Tarjan, 1973).} The deterministic
\textsc{Select} algorithm finds the $i$-th order statistic of $n$ distinct numbers
in $\Th(n)$ time in the \emph{worst case}.
\end{formalbox}

\begin{walkbox}
Why groups of $5$, and not $3$ or $7$? Redo the recurrence with groups of $g$. The
sample has $n/g$ medians, so the pivot call is $T(n/g)$; the guaranteed discard is
about $\tfrac{1}{2}\cdot\lceil g/2\rceil$ per group, giving a surviving side of
about $(1 - \tfrac{1}{2}\cdot\tfrac{\lceil g/2\rceil + 1}{g})\,n$. For the two
children to sum to less than $n$ you need $\tfrac{1}{g} + (\text{survivor
fraction}) < 1$. Check $g = 3$: the fractions come out to $\tfrac{1}{3} + \tfrac{2}{3}
= 1$ --- exactly $1$, \emph{not} less, so the proof breaks and you do \emph{not}
get linear time. Check $g = 7$: it works, but the constant is worse and sorting
groups of $7$ is fiddlier. Five is the smallest odd group size that makes the two
fractions sum below $1$ with room to spare. That is the ``magic'' of five.
\end{walkbox}

\begin{trapbox}
A tempting ``simplification'' is to use the median of medians \emph{without} the
recursion in step 3 --- say, pick the median of a fixed-size sample. That gives a
fast \emph{expected} bound but loses the worst-case guarantee, because a small
fixed sample can be unrepresentative. The recursive median-of-medians is what makes
the pivot's quality \emph{provable for every input}. Removing the recursion removes
the theorem.
\end{trapbox}

\subsection{Two algorithms, two temperaments}

You now hold two linear selection algorithms, and they have genuinely different
personalities.

\begin{center}
\begin{tabular}{@{}lll@{}}
\toprule
 & \textbf{Randomized quickselect} & \textbf{Median of medians} \\
\midrule
Expected time      & $\Th(n)$         & $\Th(n)$ \\
Worst-case time    & $\Th(n^2)$       & $\Th(n)$ \\
Randomness         & yes              & none (deterministic) \\
Hidden constant    & small ($\approx 3.4n$ cmps) & large (often $> 20n$) \\
In practice        & usually faster   & a safety net / pivot oracle \\
\bottomrule
\end{tabular}
\end{center}

In practice, randomized quickselect is the one you reach for; its constants are
small and its bad case never actually happens. Median of medians earns its keep
when you cannot tolerate \emph{any} quadratic risk, or --- elegantly --- as the
\emph{pivot oracle inside} an otherwise-randomized quicksort, giving a sort that is
$\Th(n\log n)$ in the worst case (``introselect'' and ``introsort'' use exactly
this hybrid idea).

\begin{aestheticbox}
The deeper lesson outlasts the algorithm. Median of medians is the textbook example
of \textbf{spending a little work to guarantee good structure, then proving the
guarantee pays for itself}. The recurrence $T(n) \le T(n/5) + T(7n/10) + \Th(n)$ ---
two children summing to under $1$ --- is a pattern you will meet again whenever a
clever preprocessing step shrinks two pieces that together are smaller than the
whole. Recognizing ``the children sum to less than the parent, so it is linear'' is
a transferable instinct worth far more than the specific constant $5$.
\end{aestheticbox}

\begin{historybox}
The expected-linear quickselect is due to C.~A.~R.~Hoare (1961), the same Tony
Hoare who gave us quicksort; he called it \textsc{Find}. The deterministic
worst-case-linear \textsc{Select} is the 1973 work of Manuel Blum, Robert Floyd,
Vaughan Pratt, Ron Rivest, and Robert Tarjan --- a remarkable roster, several of
them later Turing Award winners. In CLRS (4th edition) this material is
\textbf{Chapter~9, ``Medians and Order Statistics''}: Section~9.1 covers minimum
and maximum, Section~9.2 is randomized \textsc{RandomizedSelect} and its
expected-linear analysis, and Section~9.3 is the deterministic
median-of-medians \textsc{Select}. The indicator-variable expectation argument
leans on the probability background from CLRS Chapter~5, which mirrors our
Chapter~5 here.
\end{historybox}

\begin{labbox}
Implement both algorithms and \emph{count comparisons} on the same inputs.
\begin{enumerate}
  \item Generate $10{,}000$ random permutations of $1\,..\,1001$ and select the
    median with quickselect. Plot a histogram of the comparison counts. You should
    see them cluster tightly a little above $3n$, with no truly bad runs --- visceral
    proof that the $\Th(n^2)$ worst case is a phantom in practice.
  \item Feed quickselect a \emph{sorted} array with the ``always pick the last
    element'' (non-randomized) pivot rule and watch it crawl to $\Th(n^2)$. Then
    turn the randomization back on and watch the problem vanish. This is the single
    most convincing argument for why we randomize.
  \item Run median of medians on the same inputs. Confirm it never spikes --- but
    also confirm its constant is large, so for everyday work quickselect wins on
    wall-clock time despite its scarier worst case.
\end{enumerate}
\end{labbox}

\section{Exercises}

\begin{exercise}\label{ch09:ex-rank-vs-index}
For $A = [\,12, 4, 17, 9, 1, 20, 7, 15\,]$, write the sorted order and give: the
minimum, the maximum, the $3$rd order statistic, the $6$th order statistic, and the
(lower) median. For each, state both its \emph{value} and the \emph{index} it
occupies in the original (unsorted) $A$.
\end{exercise}

\begin{exercise}\label{ch09:ex-trace}
Trace \textsc{RandomizedSelect}$([\,8, 2, 11, 6, 4, 9, 1, 7\,], 1, 8, 5)$ by hand.
At each partition, choose the pivot to be the current \emph{first} element of the
subarray (so the trace is deterministic and reproducible). Show the array after each
partition, the value $q$, the rank $k$, and which branch you take. Confirm the final
answer equals the $5$th smallest element.
\end{exercise}

\begin{exercise}\label{ch09:ex-minmax}
Flesh out the walking-question pairing trick: write pseudocode that finds both the
minimum and the maximum of $n$ numbers using at most $\tfrac{3}{2}n$ comparisons.
Prove the comparison count carefully, handling the cases $n$ even and $n$ odd
separately. (Hint: the parity of $n$ changes how you seed the running min and max.)
\end{exercise}

\begin{exercise}\label{ch09:ex-second-smallest}
You can find the minimum of $n$ numbers in $n-1$ comparisons. Show that you can find
the \emph{second}-smallest in only $n + \ceil{\log_2 n} - 2$ comparisons.
(Hint: run a knockout tournament to find the minimum; the second-smallest must be
one of the $\ceil{\log_2 n}$ elements that lost directly to the eventual minimum.)
\end{exercise}

\begin{exercise}\label{ch09:ex-worst-case-input}
Describe an input and a (non-randomized) pivot rule under which quickselect performs
$\Th(n^2)$ comparisons while searching for the median. Then explain in one or two
sentences why \textsc{RandomizedPartition} makes such an input harmless.
\end{exercise}

\begin{exercise}\label{ch09:ex-groups-of-three}
Repeat the median-of-medians analysis with groups of $\mathbf{3}$ instead of $5$.
Show that the recurrence becomes $T(n) \le T(\ceil{n/3}) + T(2n/3 + c) + \Th(n)$ for
a constant $c$, that the two child sizes sum to about $n$, and argue informally that
this fails to guarantee linear time. (This is the precise reason the algorithm uses
$5$, not $3$.)
\end{exercise}

\begin{exercise}\label{ch09:ex-weighted-median}
The \textbf{weighted (lower) median} of $n$ distinct elements $x_1,\dots,x_n$ with
positive weights $w_1,\dots,w_n$ summing to $1$ is the element $x_k$ such that the
total weight of elements smaller than $x_k$ is less than $1/2$ and the total weight
of elements larger than $x_k$ is at most $1/2$. Design a $\Th(n)$ worst-case
algorithm to find it. (Hint: partition around the median-of-medians pivot, but base
the recurse-left-or-right decision on accumulated \emph{weight}, not count.)
\end{exercise}

\begin{exercise}\label{ch09:ex-quantile}
Generalize selection to the ``quantiles'' problem: given $A$ and an integer
$q \ge 2$, output the $q-1$ order statistics that divide the sorted data into $q$
groups of as-equal-as-possible size. Give an algorithm running in
$\Th(n \log q)$ time and justify the bound. (Hint: select the overall median to
split the work, then recurse on each half for the quantiles that fall in it; write
and solve the resulting recurrence.)
\end{exercise}

\begin{exercise}\label{ch09:ex-kclosest}
Given $n$ distinct numbers, a target value $t$, and an integer $k$, give an
$\Th(n)$-time algorithm that outputs the $k$ numbers closest to $t$ (ties broken
arbitrarily). The output need not be sorted. (Hint: ``closeness to $t$'' is just
$\abs{x - t}$; selection works on any key, not only the values themselves.)
\end{exercise}

\begin{exercise}\label{ch09:ex-median-two-arrays}
Two sorted arrays $X[1\,..\,n]$ and $Y[1\,..\,n]$ together hold $2n$ distinct
numbers. Give an $\Th(\log n)$-time algorithm to find the median of the combined
$2n$ elements \emph{without merging them}. (Hint: compare the medians of $X$ and
$Y$; whichever is smaller, you can safely discard its lower half and the other
array's upper half, halving the problem each step.) Explain why this beats the
$\Th(n)$ selection algorithms for this special, already-sorted case.
\end{exercise}
